% ===== PRÉAMBULE CANONIQUE — à coller VERBATIM en tête de CHAQUE .tex =====
\documentclass[11pt,a4paper]{article}
\usepackage[utf8]{inputenc}\usepackage[T1]{fontenc}\usepackage{lmodern}
\usepackage[french]{babel}
\usepackage{amsmath,amssymb,mathtools}\usepackage{icomma}
\usepackage[version=4]{mhchem}          % \ce{...} : équations & formules
\usepackage{siunitx}\sisetup{locale=FR} % \qty{}{}, \si{}, \ang{}
\usepackage{chemfig}                    % \chemfig{...} : structures développées
\usepackage{xcolor}\usepackage{colortbl}
\usepackage{tikz}\usepackage{pgfplots}\pgfplotsset{compat=1.18}
\usepackage[most]{tcolorbox}\usepackage{enumitem}\usepackage{array,booktabs}
\usepackage[a4paper,margin=2.2cm]{geometry}
\usetikzlibrary{arrows.meta,calc,angles,quotes,positioning,patterns,backgrounds,shapes.geometric}
\definecolor{primary}{RGB}{222,49,99}\definecolor{primarydark}{RGB}{150,22,55}
\definecolor{defcol}{RGB}{0,114,178}\definecolor{methcol}{RGB}{52,92,148}
\definecolor{excol}{RGB}{0,135,90}\definecolor{attncol}{RGB}{200,40,55}
\tcbset{boxbase/.style={enhanced,breakable,boxrule=0.9pt,arc=2.5pt,
  left=3mm,right=3mm,top=2.3mm,bottom=2.3mm,fonttitle=\bfseries,coltitle=white}}
% --- boîtes TUTORIEL (noms EXACTS, ne pas renommer) ---
\newtcolorbox{cle}[1][]{boxbase,colback=defcol!6,colframe=defcol,
  title={\textsc{À retenir}\ifx\relax#1\relax\relax\else\textmd{ : \itshape #1}\fi}}
\newtcolorbox{methode}[1][]{boxbase,colback=methcol!7,colframe=methcol,sharp corners,
  title={\textsc{Méthode}\ifx\relax#1\relax\relax\else\textmd{ --- \itshape #1}\fi}}
\newtcolorbox{exemplebox}[1][]{boxbase,colback=excol!5,colframe=excol,
  title={\textsc{Exemple}\ifx\relax#1\relax\relax\else\textmd{ : \itshape #1}\fi}}
\newtcolorbox{piege}[1][]{boxbase,colback=attncol!6,colframe=attncol,
  title={\textsc{Piège}\ifx\relax#1\relax\relax\else\textmd{ : \itshape #1}\fi}}
\newtcolorbox{remarque}[1][]{boxbase,colback=black!4,colframe=black!45,
  title={\textsc{Remarque}\ifx\relax#1\relax\relax\else\textmd{ : \itshape #1}\fi}}
% --- boîtes EXERCICE / CORRIGÉ (noms EXACTS, auto-comptées) ---
\newtcolorbox[auto counter]{exo}[1][]{boxbase,colback=primary!4,colframe=primary,
  title={\textsc{Exercice}~\thetcbcounter\ifx\relax#1\relax\relax\else\textmd{ --- \itshape #1}\fi}}
\newtcolorbox[auto counter]{corrige}[1][]{boxbase,colback=excol!4,colframe=excol,
  title={\textsc{Corrigé}~\thetcbcounter\ifx\relax#1\relax\relax\else\textmd{ --- \itshape #1}\fi}}
\newcommand{\R}{\mathbb{R}}\newcommand{\N}{\mathbb{N}}\newcommand{\Z}{\mathbb{Z}}\newcommand{\ds}{\displaystyle}
% (un \usepackage{fancyhdr}+en-tête est possible ; il est ignoré côté web)
% =========================================================================
\begin{document}

\begin{exo}[nitrate d'étain : lire la valence variable]
L'étain existe aux degrés $+\mathrm{II}$ et $+\mathrm{IV}$ ; le nitrate est \ce{NO3-}.
\begin{enumerate}[label=\alph*)]
  \item Écris la formule du nitrate d'étain(IV).
  \item Écris la formule du nitrate d'étain(II).
  \item Un camarade propose \ce{Sr(NO3)2} pour la question a). Explique pourquoi c'est faux.
\end{enumerate}
\end{exo}

\begin{exo}[une substance pure à partir de deux ions]
On assemble un ion d'aluminium \ce{Al^3+} avec un anion soufré.
\begin{enumerate}[label=\alph*)]
  \item Écris la formule brute de la substance formée avec l'ion \emph{sulfite}.
  \item Écris la formule brute de la substance formée avec l'ion \emph{sulfate}.
  \item Qu'est-ce qui distingue chimiquement le sulfite du sulfate ?
\end{enumerate}
\end{exo}

\begin{exo}[thiosulfate ou sulfite d'ammonium ?]
Le cation est l'ammonium \ce{NH4+}.
\begin{enumerate}[label=\alph*)]
  \item Écris la formule du thiosulfate d'ammonium.
  \item Écris la formule du sulfite d'ammonium.
  \item En quoi l'ion thiosulfate \ce{S2O3^2-} diffère-t-il de l'ion sulfite \ce{SO3^2-} ?
\end{enumerate}
\end{exo}

\begin{exo}[identifier un métal alcalin par la masse molaire]
Le phosphate d'un métal \emph{alcalin} \ce{M} a pour formule générale \ce{M3PO4}. En dissolvant
\qty{0,20}{\mol} de ce sel, on pèse \qty{42,4}{\gram}. \textbf{Données :}
$M(\ce{P})=\qty{31}{\gram\per\mol}$, $M(\ce{O})=\qty{16}{\gram\per\mol}$.
\begin{enumerate}[label=\alph*)]
  \item Justifie pourquoi la formule est bien \ce{M3PO4} (et pas \ce{MPO4} ou \ce{M2PO4}).
  \item Calcule la masse molaire du sel.
  \item Déduis-en la masse molaire de \ce{M}, puis identifie le métal.
\end{enumerate}
\end{exo}

\begin{exo}[retrouver un hydroxyde par sa composition massique]
Un composé de formule \ce{M_{$x$}(OH)_{$y$}} contient, en masse, $57{,}5\,\%$ de l'élément métallique
\ce{M}, dont la masse molaire vaut \qty{23}{\gram\per\mol}. \textbf{Donnée :}
$M(\ce{OH})=\qty{17}{\gram\per\mol}$.
\begin{enumerate}[label=\alph*)]
  \item Identifie le métal \ce{M}.
  \item Sachant que \ce{M} est monovalent, donne $x$ et $y$ et écris la formule.
  \item Vérifie ton résultat en recalculant le pourcentage massique de \ce{M}.
\end{enumerate}
\end{exo}

\end{document}
